\paragraph{About this part.} Part~I sets up the problem before any model appears. It explains what a
limit order book is and how it is built from exchange messages, fixes the notation and the probability
used throughout, and makes the survey's central distinction: predicting the next price move is not the
same as making money from it. Section~\ref{sec:intro} covers the market, the notation and the three
principles that organise the survey; Section~\ref{sec:method} explains how the literature was chosen
and how strongly each kind of claim should be believed.

% ================================================================
\section{Introduction}
\label{sec:intro}
% ================================================================

This section introduces the object of study and the language used to describe it. It explains how an
electronic market matches buyers and sellers through a limit order book, what the book looks like as
data, and how exchanges publish it; then fixes the notation and the few ideas from probability that
the rest of the survey needs. It ends with the survey's main argument---prediction, execution and
simulation are different problems---and the three principles and four axes used to organise the
literature.

\subsection{The limit order book as a dynamic market state}

Modern electronic markets operate as \emph{continuous double auctions}: buyers and sellers may
submit orders at any moment, and a trade happens whenever a buy price meets a sell price. The study
of how these trading mechanisms turn orders and information into prices is called \emph{market
microstructure}~\citep{ohara1995,harris2003,hasbrouck2007}. \emph{Liquidity}---how easily one can
trade without moving the price---is supplied by resting limit orders and consumed by aggressive
orders, and the visible limit order book records the outcome at a resolution that ranges from
individual messages to aggregated price levels~\citep{gould2013,bouchaud2018}.

This survey is written for readers with a computer-science or electrical-engineering background
and no prior finance. Market, statistical and machine-learning terms are explained where they first
appear, and Appendix~\ref{app:glossary} collects them with references. Most examples use
\emph{futures}: standardised exchange-traded contracts to buy or sell an asset at a set future
date~\citep{hull}. The running examples are the E-mini S\&P~500 future (ES, on an index of 500
large US stocks), the E-mini Nasdaq-100 future (NQ) and the 10-year US Treasury note future (ZN),
all traded on CME Group's electronic platform. Each such contract is an \emph{instrument}.

\paragraph{Orders and order types.}
An \emph{order} is an instruction to buy or sell a quantity of a contract. The quantity is counted
in \emph{lots} (one futures contract, or one share), and prices are restricted to multiples of the
\emph{tick}, the smallest allowed price step. The main order types are the
following~\citep{harris2003}.
\begin{itemize}[leftmargin=1.5em,itemsep=1pt]
\item A \emph{limit order} says ``buy (or sell) up to this many lots, at this price or better''.
  Whatever cannot trade immediately waits in the book. Waiting buy orders are \emph{bids}; waiting
  sell orders are \emph{asks} (or \emph{offers}).
\item A \emph{market order} says ``buy (or sell) this many lots now, at whatever price is
  available''. It trades immediately against waiting orders and never waits itself. A limit order
  priced to trade immediately (a buy at or above the best ask) behaves the same way; both are called
  \emph{aggressive}, and the trader is said to \emph{cross the spread}. A trader whose limit order
  waits in the book is \emph{passive}.
\item A \emph{cancel} removes a waiting order, completely or partly; a \emph{modify} (or
  \emph{replace}) changes its price or quantity.
\item \emph{Time-in-force} qualifiers say how long an order may wait: a \emph{day} order waits
  until the session ends; an \emph{immediate-or-cancel} (IOC) order trades what it can at once and
  cancels the rest; a \emph{fill-or-kill} (FOK) order trades in full at once or not at all.
\item An \emph{iceberg} order shows only part of its quantity; when the visible part trades,
  more appears. A fully \emph{hidden} order shows nothing. Both trade but are partly or wholly
  invisible in the book.
\end{itemize}

\paragraph{The book is a data structure.}
The \emph{limit order book} (LOB) is the set of all waiting orders for one instrument. Each side is
naturally implemented as an ordered map from price to a queue of orders: bids sorted from highest
price down, asks from lowest price up, with a hash map from order identifier to order so that a
cancel can find its order in constant time. The \emph{best bid} is the highest bid price, the
\emph{best ask} the lowest ask price. When an aggressive buy arrives, the exchange's
\emph{matching engine} trades it against the best ask, then the next price up, and so on until the
buy is filled or its limit price is reached; trading through several price levels is called
\emph{walking the book}. Figure~\ref{fig:book} shows the structure.

\begin{figure}[ht]
\centering
\begin{tikzpicture}[
  ord/.style={draw, minimum width=9mm, minimum height=5mm, font=\scriptsize, inner sep=1pt},
  px/.style={font=\small, anchor=east}]
\node[px] at (0,2.4) {ask 5000.75};
\node[px] at (0,1.8) {ask 5000.50};
\node[px] at (0,1.2) {ask 5000.25};
\node[px] at (0,0.0) {bid 5000.00};
\node[px] at (0,-0.6) {bid 4999.75};
\node[ord, fill=red!10] at (0.7,2.4) {12};
\node[ord, fill=red!10] at (0.7,1.8) {4};\node[ord, fill=red!10] at (1.7,1.8) {9};
\node[ord, fill=red!10] at (0.7,1.2) {3};\node[ord, fill=red!10] at (1.7,1.2) {10};\node[ord, fill=red!10] at (2.7,1.2) {1};
\node[ord, fill=blue!10] at (0.7,0.0) {7};\node[ord, fill=blue!10] at (1.7,0.0) {2};\node[ord, fill=blue!10] at (2.7,0.0) {15};\node[ord, fill=blue!10] at (3.7,0.0) {5};
\node[ord, fill=blue!10] at (0.7,-0.6) {20};\node[ord, fill=blue!10] at (1.7,-0.6) {6};
\draw[dashed] (-2.2,0.6) -- (4.4,0.6) node[right, font=\scriptsize] {spread $=1$ tick};
\draw[-{Stealth}] (0.25,-1.15) -- node[below, font=\scriptsize]{front of queue trades first} (3.9,-1.15);
\end{tikzpicture}
\caption{A small book. Each box is one waiting order and the number is its size in lots. Orders at
the same price form a first-in-first-out queue, front on the left. The best bid is 5000.00 with
$7+2+15+5 = 29$ lots; the best ask is 5000.25 with $14$ lots; the spread is one tick of 0.25.}
\label{fig:book}
\end{figure}

\paragraph{Queues and priority.}
Which waiting order trades first is decided by the venue's \emph{priority rule}. Most venues use
\emph{price--time priority}: better prices trade first, and among orders at the same price the one
that arrived first trades first. Each price level is then exactly a FIFO queue. Some products use
\emph{pro-rata} allocation instead, in which an incoming order is split among the waiting orders in
proportion to their sizes; at CME, for example, most futures match first-in-first-out while SOFR
interest-rate futures use a pro-rata or allocation algorithm~\citep{cme_matching}. Under
price--time priority an order's place in the queue matters, and the rules for keeping it matter
too: on CME Globex, decreasing an order's quantity keeps its place, while increasing the quantity
or changing the price sends it to the back of the queue~\citep{cme_matching}. The volume waiting
ahead of an order is its \emph{queue position}; much of this survey is about why that number
matters.

\paragraph{What a market-data feed shows.}
Exchanges publish the book through market-data feeds at three levels of detail. \emph{Level~1}
gives only the best bid and ask and their sizes. \emph{Level~2}, or \emph{market by price} (MBP),
gives the total size at each of the best few price levels. \emph{Level~3}, or \emph{market by
order} (MBO), reports every individual order: each addition, execution, cancellation and
modification, with an order identifier. MBO data provide the information needed to
reconstruct the displayed order-level queues, and therefore to estimate queue positions, subject to the
venue's matching rules and feed semantics; hidden liquidity and feed gaps remain outside what any feed
shows.

\paragraph{Example: building the book from Nasdaq ITCH.}
Nasdaq's TotalView-ITCH~5.0 feed is a widely used MBO protocol and a clear example of how a book is
built from messages~\citep{itch50}. Every message is a small fixed-layout binary record with a
one-character type, a timestamp in nanoseconds since midnight, and, for order messages, an 8-byte
\emph{order reference number} that identifies the order for the rest of the day. The order
messages are:

\begin{center}\footnotesize
\begin{tabular}{cll}
\toprule
Type & Name & Effect on the book \\
\midrule
\texttt{A}, \texttt{F} & Add Order & new order (side, shares, price) joins the back of its queue \\
\texttt{E} & Order Executed & order traded; reduce its shares by the executed amount \\
\texttt{C} & Order Executed With Price & as \texttt{E}, at a price different from the order's own \\
\texttt{X} & Order Cancel & partial cancel; reduce its shares \\
\texttt{D} & Order Delete & order cancelled; remove it \\
\texttt{U} & Order Replace & remove the old order; add a new one with a new reference number \\
\texttt{P} & Trade (non-cross) & an execution against a non-displayed order; no book change \\
\bottomrule
\end{tabular}
\end{center}

\noindent Building the book is a short event loop over these messages:
\begin{lstlisting}[language=Python]
book   = {'B': SortedDict(), 'S': SortedDict()}  # price -> deque of orders
orders = {}                                       # reference number -> order

def on_message(m):
    if m.type in 'AF':                            # add: join the back of the queue
        o = Order(m.ref, m.side, m.price, m.shares)
        book[o.side].setdefault(o.price, deque()).append(o)
        orders[o.ref] = o
    elif m.type in 'ECX':                         # execution or partial cancel
        o = orders[m.ref]
        o.shares -= m.shares
        if o.shares == 0: remove(o)
    elif m.type == 'D':                           # full cancel
        remove(orders[m.ref])
    elif m.type == 'U':                           # replace: remove, re-add at back
        old = orders[m.old_ref]
        remove(old)
        on_message(Add(m.new_ref, old.side, m.price, m.shares))

def remove(o):
    q = book[o.side][o.price]
    q.remove(o)                                   # O(1) with a linked list
    if not q: del book[o.side][o.price]
    del orders[o.ref]
\end{lstlisting}
Three details matter later in this survey. First, an execution message names the waiting order
that traded, so the feed reveals \emph{which} order in the queue was hit---the information needed
to reconstruct queue positions (Part~IV). Second, a replace carries a new reference number, so in
the reconstructed book the order leaves its place and rejoins at the back. Third, hidden orders
produce only \texttt{P} messages and never enter the book: the book a feed lets you build is the
\emph{displayed} book. ITCH is delivered over a sequenced UDP multicast protocol (MoldUDP64), so a
receiver can detect lost packets by gaps in the sequence numbers. CME's MDP~3.0 market-by-order feed,
used by the system in Part~IV, follows the same pattern with its own binary
encoding~\citep{cme_mdp3}.

At time $t$, an $L$-level description of the book is
\begin{equation}
  \mathcal{L}_t = \left\{ p^{b}_{\ell,t},\, v^{b}_{\ell,t},\, p^{a}_{\ell,t},\, v^{a}_{\ell,t}
  \right\}_{\ell=1}^{L},
\end{equation}
where $p^{b}_{\ell,t}$ and $v^{b}_{\ell,t}$ are the price and displayed volume at the $\ell$-th best
bid level ($\ell = 1$ is the best bid), and $p^{a}_{\ell,t}$, $v^{a}_{\ell,t}$ the same on the ask side.
\inwords{this is the market-by-price view of the book: a table with $L$ rows per side, each row a
price and the total number of lots waiting there. It can be computed from the order-level book by
summing each queue. It is the state of the market, and it changes every time an order arrives,
trades or is cancelled.}
The mid-price and
spread,
\begin{equation}
  m_t = \tfrac12\left(p^{a}_{1,t}+p^{b}_{1,t}\right), \qquad s_t = p^{a}_{1,t}-p^{b}_{1,t},
\end{equation}
\inwords{the mid $m_t$ is halfway between the best price anyone will buy at and the best price
anyone will sell at; it is the usual single-number ``price'' of the market. The spread $s_t$ is the
gap between them: buying and immediately selling back costs one spread.}
These are convenient summaries that discard most of the book: the distribution of liquidity across
levels, the sequence of events that produced the current state, and the queue position of any
individual order.

Throughout, numerical examples use stylised values chosen to make the arithmetic visible. They are
not empirical estimates, and this paper reports no new experiments.

\begin{example}[Same mid, different market]
Two E-mini S\&P 500 (ES) books have the same mid (5000.125) and the same one-tick spread (tick
$=0.25$). In the first, 12 lots rest at the best ask; in the second, 900. Buying 100 lots
aggressively costs several ticks more in the first book than in the second. Mid and spread do not
distinguish them.
\end{example}

Three properties make LOB modelling unusually hard. The data are \emph{event-driven}: orders arrive,
execute and cancel at irregular times. The state is \emph{high-dimensional} even when only a few
levels are observed. And the statistical properties are \emph{non-stationary} across time of day,
assets, volatility regimes and market-structure changes. A model that performs well for one
representation, horizon or market is therefore not, by that fact, a model of order-book dynamics.

\subsection{Notation}
\label{sec:notation}

This subsection fixes the symbols used throughout, and explains what each one means for a trader
watching the book. Symbols used in only one section are defined where they first appear. To make
the symbols concrete, each group is followed by its value in one running example: the book of
Figure~\ref{fig:book}, where at time $t$ we decide to place a passive buy order for one lot at
5000.00.

\begin{description}[leftmargin=1.5em,style=nextline,itemsep=4pt]
\item[Time and events]
  $t$ is time. Book events (additions, cancellations, executions) are numbered $n = 1, 2, \dots$
  and happen at times $t_1 \le t_2 \le \cdots$. A quantity written with subscript $t$ is its value
  just after everything that happened up to and including $t$; $t-$ means ``just before $t$''.
  $H > 0$ is a horizon---how far ahead we look---and $\delta$ is the tick size.\\
  \emph{What it means:} the market is a stream of events, and we look at it as a sequence of
  snapshots taken whenever we like. $H$ is the question ``how long am I willing to wait?''\\
  \emph{Example:} $\delta = 0.25$ (one ES tick), and we look $H = 1$ s ahead.

\item[Book]
  $p^{b}_{\ell,t}$, $v^{b}_{\ell,t}$, $p^{a}_{\ell,t}$, $v^{a}_{\ell,t}$, the mid $m_t$ and the
  spread $s_t$ are as defined in Section~\ref{sec:intro}. The change in the mid over the next $H$
  seconds is $\Delta m_{t,H} = m_{t+H} - m_t$.\\
  \emph{What it means:} $\Delta m_{t,H}$ is ``how much did the price go up (or down, if negative)
  while I waited''. Almost every prediction model in the literature targets some function of it.\\
  \emph{Example:} $p^b_{1,t} = 5000.00$, $v^b_{1,t} = 29$, $p^a_{1,t} = 5000.25$, $v^a_{1,t} = 14$,
  so $m_t = 5000.125$ and $s_t = 0.25$. If one second later the best bid and ask are 4999.75 and
  5000.00, then $m_{t+H} = 4999.875$ and $\Delta m_{t,H} = -0.25$: the price fell one tick.

\item[Information and prediction]
  $X_t$ is everything observable at $t$ that a model may use: the book, recent events, and anything
  derived from them. $Y_{t+H}$ is the quantity we want to predict, which becomes known only at
  $t+H$; $\widehat{Y}_t$ is a model's prediction of it, made at $t$.\\
  \emph{What it means:} $X_t$ is the model's input vector and $\widehat{Y}_t$ its output. The rule
  is that $X_t$ may contain only what was really known at $t$---no peeking at the future.\\
  \emph{Example:} $X_t$ could be the sizes of the five best levels on each side; $Y_{t+H}$ could be
  ``did the mid rise within one second?'', and $\widehat{Y}_t = 0.6$ a model's estimate that it
  will.

\item[Actions and orders]
  $a_t$ is what we do at $t$: place an order, keep it, cancel it, or cross the spread. $d$ is the
  side of an order, $d = +1$ for a buy and $d = -1$ for a sell. For a waiting order,
  $Q^{\text{ahead}}_t$ is the displayed volume ahead of it in its queue, and $q$ stands for a
  particular value of $Q^{\text{ahead}}_t$.\\
  \emph{What it means:} $d$ lets one formula serve both buyers and sellers, since a rising price is
  good for a buyer and bad for a seller. $Q^{\text{ahead}}_t$ is ``how many people are in line in
  front of me''.\\
  \emph{Example:} $a_t$ = place a bid at 5000.00, $d = +1$. The new order joins the back of the
  5000.00 queue, so $Q^{\text{ahead}}_t = 29$.

\item[Fills]
  $\tau_F$ is how long after $t$ the order is filled ($\tau_F = \infty$ if it never is),
  $t_F = t + \tau_F$ is the clock time of the fill, and $p^{\text{fill}}$ is the price received.
  $\Fill$ is the event ``the order was filled within the horizon'', $\{\tau_F \le H\}$, and
  $\mathbb{1}_{\Fill}$ is $1$ if that happened and $0$ otherwise.\\
  \emph{What it means:} a passive order is a bet that someone will come and trade with it; these
  symbols record whether, when and at what price they did.\\
  \emph{Example:} sellers trade through the whole 5000.00 queue 0.4 s later and reach our order:
  $\tau_F = 0.4$ s, $p^{\text{fill}} = 5000.00$, and since $0.4 \le 1$, $\mathbb{1}_{\Fill} = 1$.

\item[Profit]
  $\Pi$ is the profit per lot of an order decided at $t$, in price points unless stated otherwise
  (some examples use ticks); it is defined precisely in Section~\ref{sec:passive}. A passive order
  that is never filled makes $\Pi = 0$. $h > 0$ is the \emph{mark-out horizon}: how long after a
  fill we wait before scoring it.\\
  \emph{What it means:} $\Pi$ is the number we actually care about; every other symbol is a step
  on the way to it.\\
  \emph{Example:} with $h = 1$ s, if the mid one second after our fill is 4999.875, our buy at
  5000.00 is worth $4999.875 - 5000.00 = -0.125$: we were filled, and then the price moved against
  us. This is \emph{adverse selection} (Section~\ref{sec:adverse}), and it is why a fill is not
  automatically good news.

\item[Models]
  $M$ denotes a model, $M_0$ a baseline (usually the simplest reasonable one) and $M_1$ a candidate.
  $V(M)$ is how much money-relevant value $M$ delivers \emph{out of sample} (OOS)---on data that were
  not used to build it, as opposed to \emph{in sample}---for example net profit after costs, or a Sharpe ratio (average return divided
  by the standard deviation of return~\citep{sharpe1994}). $\Delta V = V(M_1) - V(M_0)$.\\
  \emph{What it means:} $\Delta V$ answers ``is the new model worth it, compared with the simple
  one, on data it has never seen?''\\
  \emph{Example:} if a linear model earns 0.010 points per trade after costs and a neural network
  earns 0.012, then $\Delta V = 0.002$ points per trade; Principle~3 below asks whether that is
  real and whether it pays for the extra complexity.

\item[Operators]
  $\P$ is ``the probability of'', $\E$ is ``the average of'', and a vertical bar means ``given''.
  Section~\ref{sec:primer} explains these in plain terms.
\end{description}

\subsection{A short probability primer}
\label{sec:primer}

Most equations in this paper are statements about probabilities and averages. The few ideas needed
are the following; a first course in probability covers all of them.

\begin{description}[leftmargin=1.5em,style=nextline,itemsep=2pt]
\item[{$\P(A)$, ``the probability of $A$''}] A number between 0 and 1: the fraction of times event
  $A$ would happen if the situation were repeated many times. Example: $\P(\Delta m_{t,H} > 0)$ is
  the fraction of moments after which the mid is higher $H$ seconds later.
\item[{$\P(A \mid B)$, ``the probability of $A$ given $B$''}] The same fraction, but counted only
  over the cases where $B$ is true. The vertical bar is read ``given''. Example:
  $\P(\Delta m_{t,H} > 0 \mid \text{bid queue much larger than ask queue})$ asks: of all the moments
  when buyers were queued much more heavily than sellers, what fraction were followed by a rise?
  Writing $\mid X_t$ means ``given everything we can observe at time $t$''---it is what a
  prediction model computes: a function from the current inputs $X_t$ to a probability.
\item[{$\E[Z]$, ``the expected value of $Z$''}] The long-run average of a random number $Z$. If a
  trade makes $+2$ with probability $0.6$ and $-1$ with probability $0.4$, then
  $\E[Z] = 0.6 \times 2 + 0.4 \times (-1) = 0.8$.
\item[{$\E[Z \mid B]$}] The average of $Z$ over only the cases where $B$ is true.
\item[{$\mathbb{1}_A$, the indicator of $A$}] Equal to $1$ if $A$ happens and $0$ otherwise, like a
  Boolean cast to an integer. Its average is the probability: $\E[\mathbb{1}_A] = \P(A)$. A useful
  consequence, used in Section~\ref{sec:passive}, is
  $\E[Z\,\mathbb{1}_A] = \P(A)\,\E[Z \mid A]$: the average of ``$Z$ if $A$ happened, else 0'' equals
  how often $A$ happens times the average of $Z$ when it does.
\item[{$\inf\{u > 0 : \text{condition}\}$}] ``The earliest time $u$ after now at which the condition
  holds''---the time a \texttt{while} loop that polls the condition would exit.
\item[{A hat, as in $\widehat{Y}_t$}] An estimate produced by a model, as opposed to the true value.
\end{description}

\subsection{Prediction versus tradeability}
\label{sec:tradeability}

Suppose a model estimates $\P(\Delta m_{t,H} > 0 \mid X_t)$, the probability that the mid rises over
the next $H$. That
number describes directional predictability. It does not say whether a trader can monetise it. An
aggressive buy pays the offer and must overcome the half spread and fees. A passive buy avoids the
spread but may never execute, and the event that makes it execute may itself carry information
about the next price move. The economically relevant sequence is
\begin{equation}
  \begin{gathered}
  X_t \;\to\; \widehat{Y}_t \;\to\; a_t \;\to\; \text{placement} \;\to\; \text{queue position}\\
  \;\to\; \text{execution} \;\to\; \text{realised profit and loss (P\&L)}.
  \end{gathered}
  \label{eq:chain}
\end{equation}
\inwords{data go into a model; the model makes a forecast; the forecast is turned into a decision
(buy, sell, wait, cancel); the decision becomes an order; the order takes a place in a queue;
whether and when it trades depends on that place; and only then is money made or lost. A forecast
can be right and still be lost at any later arrow.}

\begin{example}[Accurate but untradeable]
\label{ex:untradeable}
In ES (tick $0.25$ points, one-tick spread), a classifier is right on the sign of the next mid
move 55\% of the time, and the typical move is half a tick ($0.125$ points). The expected move per
trade in the predicted direction is
\[
  (0.55 - 0.45)\times 0.125 = 0.0125 \text{ points},
\]
while an aggressive entry costs the half spread, $0.125$ points, before fees. The expected directional edge is
one tenth of the half-spread cost, before fees, so it does not cover even the cost of entry. The accuracy is real, and the strategy loses money. The same signal used to decide
\emph{which resting passive order to keep and which to cancel} pays no spread; whether it helps is
an execution question (Part~III), not an accuracy question.
\end{example}

A related line of work asks not only for a forecast but for a measure of \emph{when the forecast
can be trusted}. UQ-LOB~\citep{manoharan2026uqlob} attaches an uncertainty module to a pretrained
LOB encoder and outputs a \emph{calibrated} distribution---one whose stated probabilities match
observed frequencies, so that events given 70\% happen about 70\% of the time~\citep{guo2017}---with
a scalar confidence. Restricting attention to the most confident tenth (\emph{decile}) of
predictions raises directional macro-F1 by 0.11--0.15 in its regression variant (F1 is the harmonic
mean of precision and recall; \emph{macro}-F1 averages it over the classes ``up'', ``down'' and
``stationary''), across seven cryptocurrency assets at 5, 10 and 15 second horizons. Selective
prediction changes the tradeability arithmetic directly: if a confident tier moved accuracy from
55\% to 70\% on a tenth of the signals, the expected edge per trade in
Example~\ref{ex:untradeable} would rise from $0.10$ to $0.40$ of the typical move, with one tenth
as many trades. Whether that clears the half spread still depends on the size of the moves in the
confident tier. UQ-LOB itself stops at directional reliability: its authors state that they
evaluate reliability, not profitability, and that a trading simulation with realistic execution
and fees is required before any claim about returns~\citep{manoharan2026uqlob}. That is the step
this survey is about.

This is the central question of the survey:

\begin{quote}
\emph{Which information extracted from the limit order book survives the complete translation from
statistical prediction to executable trading?}
\end{quote}

It changes how models should be compared. A model that improves classification accuracy by a
fraction of a percentage point is not more useful than a linear model if the improvement vanishes
after latency and costs. A model with modest directional accuracy may still be valuable if it
improves the timing or selection of passive executions.

\subsection{Prediction, execution and simulation are different problems}
\label{sec:three-problems}

LOB research frequently mixes up three different questions. Each is written below as a
probability or an expected value (Section~\ref{sec:primer}). The big $\P$ is ``the probability
of'', $\E$ is ``the average of'', and the vertical bar is read ``given'': everything to the right of
the bar is what we already know, everything to the left is what we want to know.
\begin{align}
  \text{Prediction:} &\quad \P(Y_{t+H} \mid X_t), \\
  \text{Execution:}  &\quad \P(\Fill \mid X_t, Q^{\text{ahead}}_t, a_t)
                        \ \text{or}\ \E[\Pi \mid X_t, Q^{\text{ahead}}_t, a_t], \\
  \text{Simulation:} &\quad \P\big(\{X_u\}_{u \in (t,\,t+H]} \mid X_t\big).
\end{align}
\inwords{
\begin{itemize}[leftmargin=1.5em,itemsep=1pt]
\item \emph{Prediction}~(4): given what we can see now, how likely is each possible value of some
  future quantity $Y$---for example, ``will the mid be higher in one second?'' It is a question
  about the market, and our own orders play no part in it.
\item \emph{Execution}~(5): given what we can see now, how much volume is ahead of our order, and
  what we chose to do, how likely is our order to be filled within $H$ seconds---or, in the second
  form, how much do we expect to earn? It is a question about \emph{our} order, and its answer
  changes with our action $a_t$.
\item \emph{Simulation}~(6): given what we can see now, how likely is each possible \emph{whole
  future path} of the market over the next $H$ seconds? A simulator draws sample paths from this
  distribution. The difference from~(4) is the difference between a weather simulator, which
  produces whole plausible tomorrows, and a forecast of tomorrow's maximum temperature.
\end{itemize}}
Two different activities are both called ``simulation'' in this literature, and the survey keeps them
apart. \emph{Generative market simulation} draws new market paths from a model of~(6): Hawkes
simulators, agent-based markets and the generative networks of Section~\ref{sec:deep-gen}.
\emph{Historical replay} does not draw paths at all: it reconstructs the one path that was recorded,
event by event, and evaluates the counterfactual outcome of the researcher's own orders against it---an
estimate of~(5) on observed data. The Kaspar system of Part~IV is a replay system, not a generative
simulator.
They are related but not interchangeable. A Hawkes process can generate realistic clustered order
flow without producing directional alpha. A survival model can give accurate fill probabilities
while playing no role in price prediction. A Transformer can predict short-horizon returns without
modelling queue dynamics. Each model must be evaluated against the problem it solves.

\begin{table}[ht]
\centering\small
\caption{One event stream (a day of ES market-by-order data), three models, three evaluation
criteria. Using the wrong column is a common source of misleading comparisons.}
\label{tab:three-problems}
\begin{tabularx}{\linewidth}{lXXX}
\toprule
Model & Target & Appropriate metric & Inappropriate metric \\
\midrule
Hawkes, 4 event types & event intensity & point-process log-likelihood; time-rescaling residual test & directional accuracy \\
Survival model & $\P(\tau_F \le 1\,\mathrm{s} \mid Q^{\text{ahead}}_t = q)$ & survival likelihood; calibration by queue bucket & mid-move AUC \\
Classifier & $\P(\Delta m_{t,H} > 0 \mid X_t)$ & log loss, AUC, then net P\&L & fill rate \\
\bottomrule
\end{tabularx}
\end{table}

\subsection{Three fixed positions}

Three positions recur throughout the paper.

\begin{principle}[The market is observed, not generated]
Recorded exchange data are the ground truth. A replay simulator reconstructs the historical book
from those data; it does not invent market behaviour. The only counterfactual it produces is the
execution outcome of the researcher's own orders.
\end{principle}

\begin{principle}[Language models generate hypotheses, not data]
In LLM-assisted model discovery the language model proposes candidate models. Fitting, replay,
validation and certification are carried out outside the language model, by fixed procedures that are
reproducible given the same data, configuration and random seeds. The language model does not generate
the market data used as evidence.
\end{principle}

\begin{principle}[Occam's razor: complexity is an empirical hypothesis]
Among models that deliver the same out-of-sample economic value, prefer the simplest. Each added
layer---model capacity, specialised hardware, and the breadth of the search that produced the
model---is a hypothesis that it adds value, and must earn incremental out-of-sample economic value
that survives realistic execution assumptions and the multiple-testing burden of its own selection.
\end{principle}

\noindent Occam's razor is applied here to economic value, not to statistical fit alone: a more
complex model is not rejected for being complex, but it carries the burden of showing that its
extra value is real. Section~\ref{sec:parsimony} makes this operational.

\subsection{How the field developed}
\label{sec:history}

The modelling of order books has passed through five overlapping stages.
\begin{enumerate}[leftmargin=1.8em,itemsep=2pt]
\item \textbf{Zero-intelligence models.} Orders arrive at random, with no strategy at all, and the
  matching rules do the rest~\citep{smith2003}. The surprise was how much this explains: on eleven
  London stocks, a model with one free parameter accounts for 96\% of the variance of the bid--ask
  spread~\citep{farmer2005}. Market structure, not trader intelligence, explains much of what the book
  looks like.
\item \textbf{Stochastic queue models.} The best queues are treated as random walks driven by order
  arrivals, cancellations and executions, giving closed-form results for the time to the next price
  change and the probability of its direction~\citep{cst2010,contlarrard2013}, later with rates that
  depend on queue sizes~\citep{huang2015}.
\item \textbf{Self-exciting point processes.} Hawkes models capture the clustering of events and the
  fact that many events are triggered by earlier ones; estimated branching ratios on E-mini S\&P data
  range from about $0.7$--$0.8$ to about one~\citep{bacry2015,filimonovsornette2012,hardimanbercotbouchaud2013}.
\item \textbf{Hand-engineered signals.} Order-flow imbalance~\citep{cks2014}, queue
  imbalance~\citep{gould2016qi} and the micro-price~\citep{stoikov2018} compress the book into a few
  numbers with clear economic meaning.
\item \textbf{Learned models.} From 2017, convolutional and recurrent networks~\citep{tsantekidis2017,zhang2019},
  then attention and spatial-temporal Transformers~\citep{wallbridge2020,berti2025tlob}, then pretrained
  message-level encoders and large generative models~\citep{linna2025lobert,m3lob}, together with
  hybrids that put point-process structure inside networks~\citep{shi2022sdpnhp,deepqr}.
\end{enumerate}
Each stage added capacity. Whether each also added \emph{economic} value is the question this survey
keeps asking.

\subsection{Positioning and taxonomy}
\label{sec:taxonomy}

The literature has three historical traditions. The first is econometric and reduced-form
modelling of prices, order flow and impact~\citep{kyle1985,glosten1985,cks2014}. The second is
continuous-time point-process and queueing models: Hawkes processes preserve the event-driven
structure and model excitation between event types~\citep{hawkes1971,bacry2015}, while
queue-reactive and first-passage models describe displayed liquidity and its
depletion~\citep{cst2010,huang2015}. The third is machine learning, from hand-engineered features
through CNNs, recurrent networks, Transformers, state-space models and, most recently, pretrained
message-level encoders~\citep{tsantekidis2017,zhang2019,wallbridge2020,linna2025lobert}.

\paragraph{Position relative to existing surveys.} Existing surveys and books cover parts of this
ground well (Table~\ref{tab:prior-surveys}): the LOB as a trading mechanism and its empirical
literature~\citep{gould2013,tripathi2020}, machine-learning prediction from the
book~\citep{zaznov2022}, generative simulation of order books~\citep{jain2024sim}, and, more recently,
large language models and autonomous agents in financial
trading~\citep{ding2024llmtrading,dong2025llmfinance,hua2026agentic}. These works provide taxonomies of
models and applications, but generally treat LOB modelling, trading prediction, execution and queue
position, market simulation, and automated hypothesis search as separate topics. We are not aware of a
survey that brings these strands together around the question of how information extracted from the
LOB survives the complete translation from statistical prediction to executable trading. That question
is also this survey's criterion for selecting literature: a model or method is included when it bears
on one link of the chain---representation, prediction, calibration, decision, queue position,
execution, replay, economic value, model search and multiple testing, and LLM-assisted hypothesis
generation. The individual subjects are not new; the integration, and the criterion that connects
them, are. Two distinctions follow from it and recur throughout: historical replay is distinguished
from generative market simulation, and language models are treated as generators of hypotheses, not
as sources of market data or of validation.

\paragraph{Generative LOB simulation versus historical replay.} The closest neighbouring review, on
LOB simulation~\citep{jain2024sim}, classifies, by methodology, models that \emph{generate} order-book dynamics, and evaluates them by their ability
to reproduce stylised facts and price impact. That is one of two legitimate ways of using a simulator,
and they answer different questions:
\begin{itemize}[leftmargin=1.5em,itemsep=1pt]
\item \emph{generative simulation} asks whether a model can produce plausible market trajectories,
  including ones never observed, such as stress scenarios;
\item \emph{historical replay} asks what would have happened to a particular strategy's orders against
  one observed trajectory.
\end{itemize}
To decide whether an order would have executed against the historical market, the market does not
need to be generated, only reconstructed. This survey is concerned primarily with the second question;
generative models are covered in Section~\ref{sec:deep-gen} as tools for simulation and stress testing.

\paragraph{Language models: hypothesis generator, not trading agent.} The surveys of LLMs in trading
review systems in which the language model reads market information, reasons about it and produces
trading decisions that are then backtested~\citep{ding2024llmtrading,hua2026agentic}. This survey does not
review that architecture. In the loop of Part~V the language model proposes a candidate model from a
fixed vocabulary of primitives; in the worked implementation it never receives market
prices~\citep{mayeski_latss}. Data enter only when an external, reproducible research engine compiles
the candidate, fits it, evaluates it out of sample and corrects for the number of candidates tried. The
implementation described there works on daily data; for order-book strategies the design adds replay
of each candidate's orders against the recorded book, which is the role of Part~IV.

The purpose here is therefore not another catalogue of architectures, but an organisation of the
literature around the translation from information to decision. Three questions recur: what representation contains the useful information;
what model extracts it; and does the result produce incremental economic value after execution
constraints?

\begin{table}[ht]
\centering\footnotesize
\caption{Existing surveys, books and benchmark studies, and what this survey adds.}
\label{tab:prior-surveys}
\begin{tabularx}{\linewidth}{>{\raggedright\arraybackslash}p{0.26\linewidth}>{\raggedright\arraybackslash}X>{\raggedright\arraybackslash}X}
\toprule
Work & Focus & Not covered there \\
\midrule
\citet{gould2013} & empirical facts and models of LOBs & deep learning; execution value; queue position in practice \\
\citet{abergel2016}, \citet{bouchaud2018} & books on empirical properties, models and simulation of order books and impact & recent deep and generative models; model-search statistics \\
\citet{tripathi2020} & systematic review of 103 articles (1981--2018) on the LOB as a trading mechanism and its implications for market efficiency, classified by sub-area, market, method and LOB measure & deep-learning prediction (named there only as a future direction); execution value; queue position \\
\citet{jain2024sim} & review of LOB simulation models and the stylised facts used to validate them & prediction and trading value \\
\citet{ding2024llmtrading}, \citet{dong2025llmfinance}, \citet{hua2026agentic} & surveys of LLMs as trading and financial agents: architectures, inputs, workflows, benchmarks, backtests, deployment & LOB microstructure; separation of hypothesis generation from evidence \\
\citet{zaznov2022} & survey of price-change prediction from the LOB & classical point-process and queue models; execution \\
\citet{prata2024}, \citet{briola2024} & benchmark re-tests of deep models; transaction-level evaluation & classical models; queue position; search statistics \\
This survey & information-to-decision chain: classical and deep models, queue position, deterministic replay, automated model discovery, evaluation & -- \\
\bottomrule
\end{tabularx}
\end{table}

\paragraph{Contributions.} The survey offers:
\begin{enumerate}[leftmargin=1.8em,itemsep=1pt]
\item a single chain---representation, prediction, decision, queue, execution, profit---that places
  classical and learned models in one framework;
\item a treatment of queue position, fill probability and adverse selection as modelling targets in
  their own right, with the signal-value/execution-value distinction (Section~\ref{sec:tradeability-framework});
\item a description of queue-exact deterministic replay as the evaluation instrument for execution
  claims (Part~IV);
\item an account of LLM-assisted model discovery as a statistical search process (Part~V);
\item evaluation standards that combine proper scoring, economic metrics, temporal validation and
  search-aware significance (Part~VI);
\item a reader-oriented apparatus: notation, a probability primer, plain-language explanations of every
  equation, a glossary and a point-process appendix.
\end{enumerate}

\paragraph{How to read this survey.} Readers new to markets should read Part~I in full and use the
glossary. Readers interested mainly in models should read Part~II, then Section~\ref{sec:benchmarks-dl}.
Practitioners concerned with execution should read Parts~III and IV and Section~\ref{sec:parsimony}.
Readers interested in automated research should read Parts~V and VI. The appendix on point processes
supports Section~\ref{sec:hawkes}.

\subsubsection*{Four axes for locating a piece of work}

Every model or paper in this survey answers four separate questions, and we describe it by its
answer to each. Think of the four answers as the four fields of a record; a work is a point in a
four-dimensional space.

\begin{description}[leftmargin=1.5em,style=nextline,itemsep=2pt]
\item[Representation: what does the model see?] The input format: raw book snapshots (the table
  of Equation~(1) at sampled times); event or message streams (every add, cancel and trade, in
  order); order-flow aggregates such as OFI (Section~\ref{sec:ofi}); queue state (sizes and
  positions in the queues); or the state of several instruments at once.
\item[Model class: what kind of function maps input to output?] Linear and econometric models;
  queue and first-passage models; point processes such as Hawkes; state-space models; convolutional
  and recurrent networks; Transformers; pretrained encoders; or a hybrid of these.
\item[Objective: what quantity does it output?] Price direction; return; event intensity;
  volatility; queue depletion time; fill probability; execution cost; or a simulated future.
\item[Economic role: what decision is the output for?] Aggressive trading (crossing the spread);
  passive trading (resting orders); market making (quoting both sides); execution of a given
  order; risk measurement; or simulation for research. A paper that evaluates only statistical
  accuracy and names no decision has an \emph{unspecified} economic role.
\end{description}

\noindent The axes are used in three ways.

\emph{First, to decide what can be compared.} Two results are comparable only if they share an
objective and an economic role, and in addition the data fields of Section~\ref{sec:comparable}.
Representation and model class are the things a study varies; objective and role are the things
it must hold fixed. A Hawkes model scored on event-intensity likelihood and a Transformer scored on
direction accuracy are not competitors, however often they appear in the same table.

\emph{Second, to organise the paper.} Part~II moves along the model-class axis, from linear models
to pretrained encoders, and Section~\ref{sec:events} along the representation axis. Part~III moves
along the objective axis from direction to fill probability and adverse selection. Parts~IV and V
fix the economic role---passive execution---and ask what each objective is worth for it.

\emph{Third, to expose a gap.} Table~\ref{tab:axes} places the main works discussed in this survey
on the four axes. The economic-role column is where the literature is thinnest: many of the
learned models are evaluated on accuracy alone, and the decision their output would drive is left
unspecified. The rest of the paper is about filling that column.

\begin{table}[ht]
\centering\footnotesize
\caption{Selected works placed on the four axes. ``Unspecified'' means the work evaluates
statistical accuracy without stating or evaluating a trading decision. \pov{} (Shadow Passive
Percentage-of-Volume) is the model-free execution algorithm of Section~\ref{sec:shadow}.}
\label{tab:axes}
\setlength{\tabcolsep}{2pt}
\begin{tabularx}{\linewidth}{l*{4}{>{\raggedright\arraybackslash}X}}
\toprule
Work & Representation & Model class & Objective & Economic role \\
\midrule
OFI~\citep{cks2014} & order-flow aggregate & linear & contemporaneous price change & price impact (execution) \\
Book dynamics~\citep{cst2010} & queue state & queueing / first passage & direction; queue depletion & market making, execution \\
Queue-reactive~\citep{huang2015} & queue state & queueing (Markov) & event intensity & simulation \\
Hawkes models~\citep{bacry2015} & event stream & point process & event intensity & simulation, risk \\
Queue value~\citep{moallemi2017} & queue state & queueing & value of a queue position & passive trading, market making \\
DeepLOB~\citep{zhang2019} & raw snapshots & CNN + recurrent & direction & aggressive trading (simple trading simulation) \\
TransLOB~\citep{wallbridge2020} & raw snapshots & Transformer & direction & unspecified \\
Deep OFI~\citep{kolm2023} & order-flow aggregate & recurrent / feed-forward & return, several horizons & unspecified (scored by out-of-sample variance explained) \\
LOBERT~\citep{linna2025lobert} & event stream & pretrained encoder & direction; next message & unspecified \\
UQ-LOB~\citep{manoharan2026uqlob} & event stream & encoder + uncertainty module & direction with confidence & unspecified (stated by the authors) \\
\pov{}~\citep{maciejewski_shadow} & event stream & none (rule) & execution cost & passive execution \\
\bottomrule
\end{tabularx}
\end{table}

% ================================================================
\section{Survey methodology}
\label{sec:method}
% ================================================================

A survey makes many claims, drawn from many papers that were written for different purposes. This
section explains what the survey covers, how its sources were chosen, how strongly each kind of
claim should be believed, and what has to be known about two studies before their numbers can be
compared. The last part is a practical checklist for reading any paper in this field.

\subsection{Scope: which markets and which questions}

The survey covers limit order books in four kinds of market. They share the mechanism of
Section~\ref{sec:intro} but differ in ways that change what a model can learn and whether a result
carries over from one to another.

\begin{table}[ht]
\centering\small
\caption{The markets covered, and how they differ in ways that matter for LOB models.}
\label{tab:markets}
\begin{tabularx}{\linewidth}{l*{3}{>{\raggedright\arraybackslash}X}}
\toprule
Market & How it trades & Typical data & What it means for models \\
\midrule
Futures (e.g.\ ES, NQ, ZN) & Each contract trades on one exchange, in nearly continuous sessions &
  Exchange MBO feeds such as CME MDP~3.0 & One book holds all the liquidity, so the book a model
  sees is the whole market. \\
US equities & The same stock trades on many exchanges and off-exchange venues at once & Exchange
  feeds such as ITCH, one venue at a time & One venue's book is only part of the market; flow
  elsewhere moves prices too. \\
Cryptocurrencies & Many independent exchanges, open 24 hours a day, every day & Public exchange
  APIs, often snapshots plus updates & Data are easy to obtain, but fees, participants and trading
  hours differ from traditional markets. \\
Fixed income & Government bond futures on futures exchanges; cash bonds on electronic
  inter-dealer platforms & Varies by venue & Treasury futures are large-tick instruments: the price rarely
  moves by more than a tick, queues are long, and queue position matters. \\
\bottomrule
\end{tabularx}
\end{table}

Within these markets the survey includes work whose main object is one of four things:
\emph{prediction} (forecasting prices or events), \emph{signal extraction} (building features from
the book), \emph{execution} (deciding how to place, keep and cancel orders), or \emph{simulation}
(reproducing the market for research). Two adjacent areas are deliberately left out. Equilibrium
models, which derive prices from the optimal behaviour of rational traders, and agent-based models,
which simulate many rule-following traders, appear only where they help explain an LOB model; they
are not fitted to the data the rest of the paper is about. Synthetic market data produced by
learned generative models are discussed in Section~\ref{sec:three-uses}, but, following Principle~1,
they play no part in the methodology the survey recommends.

\subsection{How the sources were chosen}

This is a \emph{narrative} survey, not a \emph{systematic review}. A systematic review fixes a
search query and a list of databases in advance, screens every result against written criteria, and
reports how many papers were kept at each step, so another team could repeat the search exactly. A
narrative survey instead selects work that illustrates and tests an argument. That is appropriate
here, because the aim is to organise the field around one question---what survives the translation
from prediction to execution---rather than to count papers. It also has a known weakness: the
selection may favour work that fits the argument. To limit that, the survey includes work that cuts
against its own position where such work exists, for example evidence that deep models given
suitable inputs reach state-of-the-art predictive accuracy (Section~\ref{sec:deep}).

A work was included if it made at least one of the following contributions: a representation of the
order book; a predictive model; a queue or execution model; an evaluation method; a public data set
or benchmark; or a design for simulation or research infrastructure. Preference was given to
peer-reviewed work and to widely cited preprints; recent preprints are included when they are the
only source for an idea, and are identified as preprints in the bibliography. Every reference was
checked against the publisher, the preprint server or the official documentation.

\subsection{How strongly to believe a claim}

Not all statements in a survey deserve the same confidence. We distinguish four kinds, and the
reader should ask of every claim which kind it is.

\begin{description}[leftmargin=1.5em,style=nextline,itemsep=3pt]
\item[Established result] Supported directly by a cited study, or true by mathematics.
  \emph{Example:} a linear Hawkes process is stationary if and only if the cascade of events it
  triggers dies out (Section~\ref{sec:hawkes})
  (Section~\ref{sec:hawkes}). This is a theorem; no data can overturn it.
\item[Empirical synthesis] A conclusion drawn by comparing several studies, none of which tested it
  directly. \emph{Example:} ``the input representation can matter as much as the architecture''
  (Section~\ref{sec:deep}). It rests on a pattern across papers, so it can be wrong for a market
  those papers did not study.
\item[Methodological recommendation] A proposed standard, argued for rather than proved.
  \emph{Example:} every result should report how many model configurations were tried
  (Section~\ref{sec:search}). Its justification is a statistical argument, not a measurement.
\item[Open hypothesis] Plausible, but the evidence is incomplete or conflicting. \emph{Example:}
  adding Hawkes intensities to a large network that already sees the same events adds little value
  (Section~\ref{sec:hybrid}). The survey argues for it; nobody has shown it.
\end{description}

\noindent The distinction matters because LOB studies differ in assets, horizons, targets and data,
and a result established under one set of conditions is at best a hypothesis under another.
Table~\ref{tab:evidence} in Section~\ref{sec:conclusion} classifies the main claims of this paper
in these terms.

\subsection{What makes two LOB studies comparable?}
\label{sec:comparable}

Papers in this field routinely report a single headline number---an accuracy, an F1 score, a
Sharpe ratio---and readers routinely compare those numbers across papers. Such comparisons are
meaningful only when the studies agree on the conditions that produced the numbers.
Table~\ref{tab:comparable} lists the conditions that matter most and explains how each one can
change the headline number on its own, with no change to the model.

\begin{table}[ht]
\centering\small
\caption{Fields a study must report before its results can be compared with another's.}
\label{tab:comparable}
\begin{tabularx}{\linewidth}{>{\raggedright\arraybackslash}p{0.27\linewidth}>{\raggedright\arraybackslash}X}
\toprule
Field & Why it changes the result \\
\midrule
Asset, exchange, period, trading hours & Markets differ (Table~\ref{tab:markets}); the same market
  differs between calm and volatile years and between the open, midday and the close. \\
Tick size, typical spread, depth & A large tick makes mid changes rarer and queues longer, which
  changes both what is predictable and what is tradeable. \\
Message rate & A busy instrument produces many events per second; a horizon of ten events is a few
  milliseconds in one market and minutes in another. \\
Data level (L1, L2, L3) & Only L3 data reveal the causes of book changes and queue positions
  (Section~\ref{sec:events}). \\
Sampling clock & An \emph{event clock} samples after a fixed number of events, a \emph{calendar clock}
  after a fixed number of seconds. The event clock over-represents busy periods, which behave differently. \\
Horizon and target & Predicting the sign of the next mid change is a different problem from
  predicting a move larger than a threshold, or a return over ten seconds
  (Section~\ref{sec:deep}). \\
Class balance & If most samples belong to one class, a model that always predicts that class
  scores high accuracy while knowing nothing (Example~\ref{ex:imbalance}). \\
Train/test protocol & A random split lets the model train on data from after its test period; a
  chronological split does not (Section~\ref{sec:validation}). \\
Costs, latency and execution model & The same forecast is profitable or unprofitable depending on
  whether the spread, fees and delays are counted (Section~\ref{sec:tradeability}). \\
Number of configurations tried & The best of many tries looks good by chance
  (Section~\ref{sec:search}). \\
\bottomrule
\end{tabularx}
\end{table}

\begin{example}[Accuracy without knowledge]
\label{ex:imbalance}
A three-class target labels each sample ``up'', ``down'' or ``stationary'' according to whether
the mid moves by more than a threshold within the horizon. Suppose that, with a short horizon on a
large-tick future, 80\% of samples are ``stationary'', because the mid rarely moves at all. A model
that ignores its input and always answers ``stationary'' is then 80\% accurate. A genuinely useful
model might be 75\% accurate and still be better, because it is right on some of the up and down
moves that are the only cases a trader cares about. Accuracy alone cannot show this; balanced
accuracy, macro-F1 or the confusion matrix (the table of predicted against actual classes) can.
\end{example}

\begin{example}[Two ``65\% accuracy'' results]
One study reports 65\% on FI-2010~\citep{ntakaris2018}: five Nasdaq Nordic stocks, ten trading
days, a horizon of $10$ in the data set's own sampling units, three classes. Another reports 65\% on a cryptocurrency pair sampled
on a one-second calendar clock with two classes. With two balanced classes, 50\% is what guessing
achieves; with three classes, guessing achieves anywhere from 33\% to the share of the largest
class. The two numbers sit on different scales and share almost none of the fields of
Table~\ref{tab:comparable}. Their equality carries no comparative information.
\end{example}

\subsection{A reader's checklist}

The following questions can be asked of any paper in this field, including this one. A study that
cannot answer them has not yet shown that its result is useful for trading.

\begin{enumerate}[leftmargin=1.8em,itemsep=2pt]
\item \textbf{What exactly is predicted, over what horizon, on what clock?} (Section~\ref{sec:deep})
\item \textbf{What does a trivial model score on the same data?} Always guessing the most common
  class, or a linear model on order-flow imbalance (Section~\ref{sec:ofi}).
\item \textbf{Was the test period strictly after the training period?} (Section~\ref{sec:validation})
\item \textbf{How many models or settings were tried before this one was reported?}
  (Section~\ref{sec:search})
\item \textbf{What decision would the prediction drive, and was that decision evaluated after
  spread, fees and latency?} (Sections~\ref{sec:tradeability} and~\ref{sec:passive})
\item \textbf{If the decision uses passive orders, how were fills decided?} A fill rule that ignores
  the queue overstates results (Section~\ref{sec:replay}).
\end{enumerate}

